\documentclass[10pt]{article}
\usepackage[margin=20mm]{geometry}
\usepackage{amsmath,amssymb,mathtools,fontspec,unicode-math,xurl,hyperref}
\IfFontExistsTF{DejaVu Serif}{\setmainfont{DejaVu Serif}}{\setmainfont{Latin Modern Roman}}
\setmathfont{latinmodern-math.otf}
\setlength{\parindent}{0pt}\setlength{\parskip}{6pt}\sloppy
\title{A verification of selected arguments in
Non-vanishing of Geodesic Periods of Automorphic Forms}
\author{}
\date{}
\begin{document}
\maketitle
Abstract: The quantitative edge count, graph measure, and distribution-transfer arguments in the acquired manuscript are examined. A one-dimensional Gaussian estimate in the proof of Theorem 5.1 is incorrect; a corrected calculation still recovers the headline exact-zero exponent. The graph transform loses mass at isolated vertices, which occur in the actual modular graph, and the general uniform transfer lemma needs more than continuity of its density. An Eisenstein normalization discrepancy remains unresolved. These findings do not constitute a counterexample to every stated asymptotic theorem.\par
\section*{Identity, scope and principal distinctions}
The source is Petru Constantinescu and Asbjørn Christian Nordentoft, Non-vanishing of geodesic periods of automorphic forms, arXiv:2404.12982v2, 29 July 2025, 38 physical pages; its title-page date is 30 July 2025 [CN]. The assigned DOI is 10.1007/s00039-025-00715-z. This exact author manuscript is assessed; identity with the journal version is not established. Physical and printed page numbers agree.\par
Theorem 1.3 concerns zero periods of a fixed nonzero level-one weight-zero Hecke–Maaß cusp form, with geodesics ordered by length. Theorem 1.5 treats the stated cuspidal cases for groups with a cusp. Theorem 1.9 concerns a graph-weighted distribution, whereas unweighted normality is Conjecture 1.10. The arithmetic application orders real quadratic fields by fundamental unit, not by discriminant, and asserts positive proportion rather than Michel’s question for every sufficiently large field. These distinctions are retained throughout.\par
The checks below cover selected mathematical arguments and actual cited source-page images, not an independent proof of every external input. In particular, the period-comparison proposition, the cited central limit theorems, and sparse equidistribution are used with their stated hypotheses and remain external dependencies.\par
\section*{Problem 1: a squared small-interval estimate}
Assessment: Major, quantitative proof gap. Theorem 5.1 (p. 25) claims, for \(\delta>0\), \[\#\{y\in Y_N^*:|P_f(y)|\le(\log N)^{1/2-\delta}\}
 \ll \frac{N^2}{(\log N)^{1+\min(\delta,1/4)}}.\] Its proof sets \(\alpha=1+\min(\delta,1/4)\). Equation (5.2) bounds the mass of the real standard Gaussian on \([-r,r]\), \(r=c''(\log N)^{-(\alpha-1)}\), by \(O((\log N)^{-2(\alpha-1)})\). The actual one-dimensional density satisfies \[\int_{-r}^{r}\frac{e^{-x^2/2}}{\sqrt{2\pi}}\,dx
 =\sqrt{\frac2\pi}\,r+O(r^3).\] Since \(\alpha>1\), its ratio to the proposed squared-radius bound is unbounded. For example, at \(\delta=1/8\) the mass has order \((\log N)^{-1/8}\), not \(O((\log N)^{-1/4})\). The error \(O((\log N)^{-1/2})\) in Theorem 3.4 cannot repair this mismatch. Both the printed integral and the real one-dimensional input were checked in the PDF (pp. 17, 25). A two-dimensional Gaussian estimate does not apply to these Maaß periods.\par
\subsection*{The available quantitative bound}
Put \(L=\log N\), \(\rho=\min(\delta,1/2)\), and let \(B_N\) be the neighbours of the exceptional geodesics \(A_N\). Proposition 3.1 gives \[|L_f(x)|\le L^{1/2-\delta}+O(1)
           =O(L^{1/2-\rho}),\qquad x\in B_N.\] The constant error must be retained when \(\delta>1/2\). The real Gaussian input yields \(|B_N|\ll N^2L^{-\rho}\).\par
Only a degree upper bound is needed. The trace progression gives \(\deg(x)\le N/c(x)+2\). Together with \(\#\{x:c(x)\le t\}=O(t^2)\), partial summation gives \[\sum_{x\in B_N}\deg(x)
 \ll Nt+(N/t+1)|B_N|,\qquad 1\le t\le N.\] Choose \(t=N/L^{\rho/2}\). The right side is \(O(N^2L^{-\rho/2})\). If \(|A_N|\) is already below the resulting target, nothing remains; otherwise its total length meets the threshold of Theorem 4.6. Proposition 4.2 and that sparse equidistribution input then give \(\sum_{y\in A_N}\deg(y)\gg |A_N|L\). Consequently the checked argument proves \[|A_N|\ll
 \frac{N^2}{L^{\,1+\min(\delta/2,1/4)}}.\] This is weaker than the printed assertion for \(0<\delta<1/2\). A new argument or a weaker statement is required. This is a demonstrated failure of the given quantitative derivation, not an exhibited Maaß form contradicting the theorem.\par
\subsection*{Why the headline exact-zero exponent survives}
For \(P_f(y)=0\), the comparison gives \(|L_f(x)|=O(1)\) directly. Thus \(|B_N|\ll N^2L^{-1/2}\), and the cutoff \(t=N/L^{1/4}\) gives \(|A_N^0|\ll N^2/L^{5/4}\). Since \(|\operatorname{tr}\gamma|=2\cosh(\ell(C_\gamma)/2)\), dyadic summation recovers \(O(e^X/X^{5/4})\). This is a verified repair of the exact-zero deduction conditional on the named source inputs. It is counterevidence to a claim that Problem 1 invalidates the headline non-vanishing result.\par
The qualitative cuspidal edge argument also needs only small/large normalized-period probabilities tending to zero. Weak Gaussian convergence, absence of an atom at zero, and tightness provide that conclusion by comparison with fixed intervals or rectangles. The degree cutoff can then grow slowly enough to make the exceptional edge mass \(o(N^2)\). The geometric lower bound for a positive-proportion exceptional collection contradicts this. Isolated vertices have degree zero and do not affect this particular edge-sum argument.\par
\section*{Problem 2: the graph transform need not have mass one}
Assessment: Major, missing condition in the weighted-transfer route. For the definition (6.1), p. 28, let \(I=\{x\in X:\deg(x)=0\}\). Direct summation gives \[G(\mu)(Y)
 =\sum_{\deg(x)>0}\sum_{y\in e(x)}\frac{\mu(x)}{\deg(x)}
 =1-\mu(I).\] Thus this is a probability measure only if \(\mu(I)=0\). The definition \(e^{-1}(B)=\{x:e(x)\subset B\}\) on p. 18 includes isolated vertices. Taking \(B=Y\) in Lemma 6.1 would therefore give \(1\le 1-\mu(I)\), which is false when they carry positive mass. The observation that a vertex adjacent to a given \(y\) has positive degree does not exclude isolated vertices elsewhere.\par
This issue occurs in the actual modular graph, not merely in an unrelated abstract graph. Take \[N=21,\qquad
 \gamma=\begin{pmatrix}-6&-5\\17&14\end{pmatrix}.\] Its determinant is one, \(c_\gamma=17\le21\), and its trace is \(8\). All representatives of its double coset have traces \(8+17k\), \(k\in\mathbb Z\). The closest values to zero are \(8\) and \(-9\); all others have absolute value at least \(25\). None belongs to the required absolute-trace band \([21/2,21]\). Hence its vertex in \(X_{21}\) has degree zero and positive uniform mass.\par
There are also arbitrarily large examples. For every integer \(m\ge12\), take \[\gamma_m=\begin{pmatrix}m+1&1\\2m+1&2\end{pmatrix},
 \qquad N_m=\left\lfloor\frac54(2m+1)\right\rfloor .\] The determinant is one and the traces in the double coset are \(m+3+k(2m+1)\). The two closest values, \(m+3\) and \(2-m\), have absolute value below \(N_m/2\); the next values have absolute value exceeding \(N_m\). This again gives an isolated vertex.\par
\subsection*{What is established and what remains necessary}
These examples refute exact normalization for the graph as defined. They do not determine the asymptotic proportion \(\zeta_N=\mu_N(I_N)\). In particular, they alone are not a counterexample to a Gaussian limiting assertion.\par
The general lemma can be repaired by requiring \(\mu(I)=0\), but this hypothesis cannot simply be inserted into the actual application. One possible route is to restrict and normalize \(\mu_N\) on the nonisolated vertices, and prove an adequate bound on \(\zeta_N\). The resulting distributional comparison has an additional \(O(\zeta_N)\) term when \(\zeta_N\) is small. At least \(\zeta_N=o(1)\) is necessary for the intended unit-mass limit; the effective claim needs a corresponding rate. No such bound or replacement graph construction is established in this review. The Major assessment concerns the unverified transfer route for a stated main weighted-distribution result, without asserting that its asymptotic conclusion is false.\par
\section*{Problem 3: continuity does not supply a uniform density bound}
Assessment: Minor, general auxiliary lemma. Lemma 6.2 (pp. 29–30) asserts a uniform interval error \(O_F(h(n)+E(n))\) for any continuous probability density \(F\). Its proof uses \(\int_b^{b+E}F=O_F(E)\) uniformly in \(b\). A continuous integrable density can be unbounded.\par
Here is a counterexample to the lemma itself, with no isolated vertices. Set \(w_j=2^{-j}/j\) and \[F(x)=\sum_{j\ge1}j\left(1-\frac{|x-3j|}{w_j}\right)_+ .\] The disjoint triangular spikes have masses \(jw_j=2^{-j}\), so \(F\) is a continuous probability density. For each \(n\), take a finite, equally weighted quantile approximation to \(F\) whose uniform interval error is at most \(h(n)\le w_n^2\). Such an approximation is obtained from midpoint quantiles with sufficiently many atoms.\par
Let the graph be the matching between two copies of these atoms, let \(f_n\) equal their quantile values, and put \(g_n=f_n+w_n\). Every edge has error \(E(n)=w_n\), and the graph transform is a probability measure. On the interval \([0,3n]\), the output distribution differs from the target by \[\int_{3n-w_n}^{3n}F(x)\,dx+O(h(n))
 =\frac{nw_n}{2}+O(w_n^2).\] Its ratio to \(h(n)+E(n)\) tends to infinity. All stated hypotheses of the lemma hold, so its uniform rate is false as quantified.\par
Requiring \(F\) to be bounded gives the claimed endpoint estimate and repairs this step. More generally the error should use the concentration modulus \(\sup_b\int_b^{b+O(E)}F\), which tends to zero but need not be \(O(E)\). The Gaussian densities used in the application are bounded; this is why the density defect is local to the general lemma. It does not resolve the separate isolated-mass issue in Problem 2.\par
\section*{Remaining checks and verification limits}
\subsection*{Period comparison and an Eisenstein dependency}
On an edge, \(c_\gamma\le N\) and \(|\operatorname{tr}\gamma|\ge N/2\). Therefore the cuspidal error in Proposition 3.1 is \(O_f(1)\), and after division by \(\sqrt{\log N}\) it is \(O_f((\log N)^{-1/2})\). The sign \((-1)^{k/2+1}\) is retained in §6.3 and does not change a centred real or complex Gaussian law. The external proof underlying Proposition 3.1 has not been independently reconstructed here.\par
For Eisenstein periods there is a separate unresolved normalization discrepancy. Theorem 3.5 on p. 17 visibly divides the vertical period by \(c_{E^*}\sqrt{\log N}\), while pp. 26–27 use \(\sqrt{\log N}(\log\log N)^{3/2}\). The latter matches Theorem 1.6’s length scale \(\sqrt X(\log X)^{3/2}\). The quoted external central limit theorem has not been checked to decide which printed normalization is correct. This is recorded as a NeedsEvidence dependency; the Eisenstein deduction is not certified.\par
\subsection*{The positive-proportion arithmetic mechanism}
In the stated Heegner setup of §7, character orthogonality turns Waldspurger’s identity into a sum of squared geodesic periods. If every relevant central value vanishes, all those periods vanish. Conditional on the source’s cuspidal non-vanishing and class-number inputs, \[\sum_{D\ {\rm good}}h^+(D)\gg \frac{X^2}{\log X},
 \qquad
 \sum_D h^+(D)^2\ll \frac{X^3}{(\log X)^2}.\] Cauchy–Schwarz then gives \(\#\{D\ {\rm good}\}\gg X\). Together with the counting input this yields a positive proportion. This checks the mechanism, not the printed optimal constant or all wide/narrow class-number conventions. The extension from the explicit square-free level setup of §7 to every level in Corollary 1.7 is not independently verified. No pointwise-in-discriminant conclusion follows from this averaged argument.\par
\section*{Conclusion}
There are two supported Major proof/definition gaps and one Minor auxiliary-lemma defect within the inspected scope. The exact-zero exponent and the qualitative cuspidal edge-count deduction have verified repairs or conditional supporting arguments; a blanket rejection of all non-vanishing results would exceed the evidence. The stronger quantitative small-period bound, the normalization needed for weighted transfer, and the Eisenstein scale require the specified additional work. Author adoption and journal-version changes are unknown. Scientific decision remains pending, human reviews are empty, and this initial continuation is not independent acceptance.\par
\section*{Scope and references}
The Gaussian expansion, matrix trace examples, density construction and edge optimization are analytic checks, not numerical experiments. The original source extraction is reused with its historical execution origin. No author code was executed. Unavailable optional code is not a mathematical defect. External theorems are declared inputs, not relabelled as freshly verified proofs.\par
1 P. Constantinescu and A. C. Nordentoft, Non-vanishing of geodesic periods of automorphic forms, arXiv:2404.12982v2, 29 July 2025.\par
\end{document}
